\begin{pastexam}[label={px:nccu-108-dsa-02}]{政大資科所 108 學年度 第 2 題}
（問答）

(10\%) Give two advantages of hash tables over balanced binary search trees.
\end{pastexam}

\begin{pastsoln}
下面 $n$ 是存的鍵數。

\textbf{優點一：搜尋、插入、刪除的期望時間是 $O(1)$}。 雜湊表用雜湊函數直接算出鍵該放的位置
（資結書第~12.1~節），只要負載因子 $\alpha = n/m$ 維持在常數以下（例如超過時就把表加大、重新雜湊），
在簡單均勻雜湊的假設下，鏈結法搜尋不在表裡的鍵期望比較 $\alpha$ 個鍵、在表裡的鍵不到
$1 + \alpha/2$ 個（資結書定理~12.2），加上算雜湊值的 $O(1)$，是常數，和 $n$ 無關。 平衡二元搜尋樹
（AVL、紅黑樹）每次操作都要從根往下比較，$\Theta(\log n)$。 $n = 10^6$ 時樹要比大約 $20$ 到 $40$ 次，
雜湊表平均一兩次。

\textbf{優點二：鍵不需要能比大小}。 二元搜尋樹靠「比根小往左、比根大往右」，鍵一定要有全序。 雜湊表
只需要一個雜湊函數和「相等」的判斷，例如字串、座標、檔案內容這類不方便定義大小、或比大小很花時間的
鍵，都可以直接放。 另外，雜湊表的實作比平衡樹簡單：不需要旋轉、顏色或高度的維護，程式短、比較不容易
寫錯。

（取捨：雜湊表最壞情況是 $\Theta(n)$，例如所有鍵撞在同一格；而且它不保存鍵的順序，找最小值、
找前一個或後一個、列出某個範圍的鍵都要 $\Theta(n)$，這些是平衡樹的強項，見資結書第~10.5~節的比較。）
\end{pastsoln}
